\documentclass[10pt,t,english,aspectratio=1610]{beamer}
\usepackage{fontawesome}
\usepackage{graphicx}
\usepackage{array}
\usepackage[normalem]{ulem}
\usepackage{amsfonts,amsmath,amssymb,bm,bbm}
\usepackage{mathrsfs}
% NOTE: sgame is intentionally NOT used -- its `game` environment is
% incompatible with the `array` package (required by colortbl/makecell/multirow)
% and errors with "Missing # inserted in alignment preamble". We use the modern
% TikZ-based `xgames` package instead (loaded after TikZ below).
\usepackage{graphicx,pstricks}
\usepackage{xcolor}
\usepackage{colortbl}
\usepackage{makecell}
\usepackage{tikz,tikzsymbols,gnuplottex}
\usetikzlibrary{decorations.pathreplacing,shapes,positioning,arrows.meta,fit,backgrounds,calc}

% --- xgames: TikZ-based game-theory typesetting (matrixgame, gametree,
% beliefspace, automaton). Not on CTAN; xgames.sty + fikz.sty are vendored at
% the repo root. CC BY-SA 4.0 (xgames) / GPL-3.0 (fikz), (c) Benjamin Bernard.
\usepackage{xgames}

% Custom command for shading rows in a matrixgame
% ---------------------------------------------------------------------------
% \shaderow[<fill options>]{<row>}{<first column>}{<last column>}
% Call this inside the matrixgame, after \payoffmatrix
% Can wrap in a \foreach to shade multiple regions
% Example: \foreach \i in {1,...,5}{ \shaderow[green!50!black]{\i}{\i}{5} }
\newcommand{\shaderow}[4][red!60]{%
  \ifnum#3=1\relax\def\xl{0}\else\pgfmathsetmacro{\xl}{\pmtwi[#3-1]}\fi
  \ifnum#2=1\relax\def\yt{0}\else\pgfmathsetmacro{\yt}{\pmthi[#2-1]}\fi
  \pgfmathsetmacro{\xr}{\pmtwi[#4]}%
  \pgfmathsetmacro{\yb}{\pmthi[#2]}%
  \fill[opacity=0.18, #1] ($(pmt1-tl)+(\xl,-\yt)$) rectangle ($(pmt1-tl)+(\xr,-\yb)$);%
}

\usepackage[english]{babel}
\usepackage[utf8]{inputenc}
\usepackage{appendixnumberbeamer}
\usepackage{datetime2}
\usepackage{booktabs}
\usepackage{setspace}
\usepackage{rotating}
\usepackage{multirow}
\usepackage{listings}

\usepackage{natbib}
\bibliographystyle{unsrtnat}
\renewcommand{\bibfont}{\footnotesize}
\renewcommand{\bibsection}{\vspace{-1em}}

%\usepackage{matlab-prettifier} % For enhanced MATLAB highlighting
\lstset{
    language=matlab,
    % Other options for styling (optional)
    basicstyle=\ttfamily\footnotesize, % Font style
    keywordstyle=\color{blue}, % Keyword color
    commentstyle=\color{green!50!black}, % Comment color
    stringstyle=\color{red!70!black}, % String color
    numbers=left, % Line numbers on the left
    numberstyle=\tiny\color{gray}, % Style for line numbers
    frame=single, % Frame around the listing
    breaklines=true, % Allow line breaking
    captionpos=b, % Caption at the bottom
    tabsize=4 % Tab size
}

% ------------------------------------------------------------------------------
% Use the beautiful metropolis beamer template
% ------------------------------------------------------------------------------
\usepackage[T1]{fontenc}
\usepackage[utf8]{inputenc}
\usepackage{fontawesome}
\usepackage{FiraSans} 
\mode<presentation>
{
  \usetheme[progressbar=foot,background=light]{metropolis} 
  \usecolortheme{default} % or try albatross, beaver, crane, ...
  \usefonttheme{default}  % or try default, serif, structurebold, ...
  \setbeamertemplate{navigation symbols}{}
  \setbeamertemplate{caption}[numbered]
  %\setbeamertemplate{frame footer}{My custom footer}
  \makeatletter
    \setbeamertemplate{frametitle continuation}{}
  \makeatother
} 
\addtobeamertemplate{frametitle continuation}{}{%
  \setbeamertemplate{frametitle continuation}{}%
}
\setbeamertemplate{frametitle continuation}{}

\newenvironment{stepenumerate}{\begin{enumerate}[<+->]}{\end{enumerate}}
\newenvironment{stepitemize}{\begin{itemize}[<+->]}{\end{itemize} }
\newenvironment{stepenumeratewithalert}{\begin{enumerate}[<+-| alert@+>]}{\end{enumerate}}
\newenvironment{stepitemizewithalert}{\begin{itemize}[<+-| alert@+>]}{\end{itemize} }

\newtheorem{question}{Question}
\newtheorem{claim}{Claim}
\newtheorem{proposition}{Proposition}
\newtheorem{remark}{Remark}
\newtheorem{conjecture}{Conjecture}

\definecolor{metrop}{RGB}{29, 44, 44}
\colorlet{GrayLight}{black!15}
\colorlet{GrayMedium}{black!30}
\colorlet{ForestGreen}{green!60!black}

\newenvironment{transitionframe}{
  \setbeamercolor{background canvas}{bg=black!80}
  \begin{frame}}{
    \end{frame}
}

\newcommand{\bsk}{

\bigskip

}

\newcommand{\pd}{\partial}
\newcommand{\RR}{\mathbb{R}}

\newcommand*\hugme[1]{\tikz[baseline=(char.base)]{\node[shape=ellipse,draw,inner sep=0pt] (char) {#1};}}

\newcounter{saveenumi}
\newcommand{\seti}{\setcounter{saveenumi}{\value{enumi}}}
\newcommand{\conti}{\setcounter{enumi}{\value{saveenumi}}}
\resetcounteronoverlays{saveenumi}



\newcommand\dotprod[2]{\langle #1 , #2 \rangle}
\newcommand{\ft}[1]{\widehat #1}
\newcommand{\qabove}[1]{\overset{\text{\large \textbf ?}}{#1}}
\newcommand{\eqae}{\overset{\text{a.e.}}{=}}
\newcommand{\calp}{\mathcal{P}}
\newcommand{\calg}{\mathcal{G}}
\newcommand{\calb}{\mathcal{B}}
\newcommand{\textd}{\text{d}}
\newcommand{\bbr}{\mathbb{R}}
\newcommand{\binm}{\mathbin{M}}
\newcommand{\binc}{\mathbin{C}}
\newcommand{\binb}{\mathbin{B}}
\newcommand{\calc}{\mathcal{C}}
\newcommand{\calh}{\mathcal{H}}
\newcommand{\bfone}{\mathbf{1}}
\newcommand{\bbe}{\mathbb{E}}
\newcommand{\bfle}{\mathbf{e}}
\newcommand{\calf}{\mathcal{F}}
\newcommand{\cala}{\mathcal{A}}
\newcommand{\cale}{\mathcal{E}}
\newcommand{\bbn}{\mathbb{N}}
\newcommand{\cantor}{\calc}
\newcommand{\calY}{\mathcal{Y}}
\newcommand{\textb}{\text{B}}
\newcommand{\calm}{\mathcal{M}}
\newcommand{\bint}{\mathbin{T}}
\newcommand{\ep}{\epsilon}
\newcommand{\bbq}{\mathbb{Q}}
\newcommand{\bbp}{\mathbb{P}}
\newcommand{\cals}{\mathcal{S}}
\newcommand{\emptysequence}{e}
\newcommand{\bbz}{\mathbb{Z}}
\newcommand{\fraka}{\frak{A}}
\newcommand{\frakb}{\frak{B}}
\newcommand{\length}{\text{length}}
\newcommand{\bfn}{\mathbf{N}}
\newcommand{\support}{\text{support}}
\DeclareMathOperator*{\argmax}{arg\,max}
\newcommand{\dom}{\mbox{dom}}
\def\ut{\underline t}
\def\um{\underline m}
\def\PP{\mathbb{P}}
\def\EE{\mathbb{E}}
\def\RR{\mathbb{R}}

% Title page info
\title[Incomplete Information: Dynamic]{Chapter 15: Incomplete Information, Dynamic Games}
\author[ECON 5001]{ECON 5001\\Prof. Paul J. Healy\\Based on: \textit{Game Theory} by Giacomo Bonanno} \color{metrop}
\institute[OSU]{}
\date[]{\vfill {\tiny Updated \today\ at\ \DTMcurrenttime}}

% ---------------------------------------------------------------------
% SCOPE AND NOTATION
%   Builds on decks 11-13. Solution concept: WSE of the Harsanyi game.
%   Section 15.1's three examples: the small perfect-info example
%   (Figs 15.1-15.5), reputation (chain store), and signaling (strikes),
%   plus truth-in-advertising in words. Section 15.2 (multi-sided,
%   bargaining) gets one footnote.
%   Chain store is re-skinned with deck 05's IKEA/Wayfair numbers:
%   out (1,5), fight (0,0), share (2,2); a hotheaded IKEA gets 3 from
%   fighting. Bonanno uses share (1.5,1.5), hotheaded fight 2, and gets
%   p >= 1/3. With deck 05's numbers the threshold is p >= 1/2.
%   Payoffs in the chain store are (entrant, IKEA) per city.
%
% ANSWERS (all checked numerically)
%   Example (Bonanno 15.1-15.5): P1 thinks alpha (P2's payoffs: D 2,
%   A 1) with prob 1/3, beta (D 1, A 2) with prob 2/3. True state alpha.
%     Strategic form, P2 strategy = (at alpha, at beta):
%       B: (1,0) for every P2 strategy
%       T: DD (0, 4/3)  DA (4/3, 2)  AD (2/3, 1)  AA (2, 5/3)
%     BNE: (B, DD), (B, AD), (T, DA). Only (T, DA) is a SPNE / WSE.
%     True game alpha: BI is B. But (T, DA) plays T then D.
%   Chain store (deck 05 numbers): candidate WSE
%     Wayfair (Columbus): Out. IKEA (either type): fight entry in Columbus.
%     Living Spaces (Cincinnati): out after Out, in after Share,
%       out after Fight. Beliefs Pr(hotheaded): p, 0, 1.
%     After Out: in gives 2(1-p), out gives 1 -> out iff p >= 1/2.
%     Rational IKEA: fight 0 + 5 = 5 > share 2 + 2 = 4.
%     Hotheaded IKEA: fight 3 + 5 = 8 > share 2 + 3 = 5.
%     Wayfair: in -> fought -> 0 < out 1.
%   Strikes (Bonanno's numbers): piH = 100, piL = 55, wH = 60, wL = 50,
%   alpha = Pr(high) = .7, delta = .6.
%     Separating WSE: ask wH; high firm accepts, low rejects; then ask
%     wL, accepted.
%     Union in period 2 after a rejection: p = Pr(low) = 1 by Bayes.
%       Needs p >= 1 - wL/wH = 1/6.
%     High firm won't mimic: 100 - 60 = 40 >= .6(100 - 50) = 30.
%     Low firm: reject gives .6(5) = 3 > 55 - 60 = -5.
%     Union period 1: .7(60) + .3(.6)(50) = 42 + 9 = 51 >= 50.
%     True state: low profit -> a strike happens.
%   Worked example (Honors signal). Pr(Star) = 1/3. Worker: Manager job
%   5, Analyst job 2, minus honors cost (Star 1, Average 4).
%   Nationwide: Manager to Star 4, Manager to Average 0, Analyst 2.
%     Nationwide hires Manager iff mu(Star) >= 1/2.
%     Separating WSE: Star takes Honors, Average doesn't. Honors ->
%       Manager (mu = 1), No -> Analyst (mu = 0).
%       Star: 4 > 2. Average: 2 > 1.
%     Pooling WSE (nobody takes Honors): No -> Analyst (mu = 1/3).
%       After Honors (off path): Analyst, needs mu(Star | Honors) <= 1/2.
%     Reasonable? No: Average's Honors payoff is at most 1 < 2, so only
%       a Star would ever take it. Belief should be mu = 1.
%     No pooling on Honors (Star would deviate to No).
%     Average's cost 2 instead of 4: no separating WSE (Average would
%       mimic: 5 - 2 = 3 > 2). Only the pooling one survives.
% ---------------------------------------------------------------------

% The two possible games in the opening example (xgames won't take
% macro arguments inside payoff lists, so they're written out)
\newcommand{\minigamealpha}{%
\begin{gametree}[player names={1, 2}, labels=short, scale=0.85]
  \branch[player=1, root=w, label=90] from (0,0) to[v=1.2, h=2.4] {B(1,0); T};
  \branch[player=2, label=0] from (w2) to[v=1.2, h=1.8] {D(0,2); A(2,1)};
\end{gametree}}
\newcommand{\minigamebeta}{%
\begin{gametree}[player names={1, 2}, labels=short, scale=0.85]
  \branch[player=1, root=w, label=90] from (0,0) to[v=1.2, h=2.4] {B(1,0); T};
  \branch[player=2, label=0] from (w2) to[v=1.2, h=1.8] {D(0,1); A(2,2)};
\end{gametree}}

% Harsanyi transformation of the opening example (Bonanno Fig 15.4)
\newcommand{\examplegame}[1][0.72]{%
\begin{gametree}[player names={1, 2}, labels=short, scale=#1]
  \branch[player=0, root=w, label=90] from (0,0) to[v=1.2, h=5.8] {$\alpha$\ $\frac13$; $\beta$\ $\frac23$};
  \branch[player=1, label=none] from (w1) to[v=1.3, h=2.2] {B(1,0); T};
  \branch[label=none] from (w2) to {T; B(1,0)};
  \branch[player=2, label=180] from (w12) to[v=1.3, h=1.5] {D(0,2); A(2,1)};
  \branch[player=2, label=0] from (w21) to {D(0,1); A(2,2)};
  \informset[player=1] (w1) -- (w2);
\end{gametree}}

% Chain store knowledge-belief structure
\newcommand{\chainks}{%
\begin{tikzpicture}[x=1.5in, y=0.5in]
  \node at (0, 0.8) {$\alpha$: rational}; \node at (1, 0.8) {$\beta$: hotheaded};
  \foreach \y/\n/\c in {0/IKEA/player2, -1/Wayfair/player1, -2/Living Spaces/player3}
    {\node[anchor=east, text=\c] at (-0.3, \y) {\n:};
     \fill (0, \y) circle (2pt); \fill (1, \y) circle (2pt);}
  \foreach \y/\c in {0/player2, -1/player1}
    {\draw[\c, thick, rounded corners=7pt] (-0.15, \y-0.25) rectangle (0.15, \y+0.25);
     \draw[\c, thick, rounded corners=7pt] (0.85, \y-0.25) rectangle (1.15, \y+0.25);}
  \draw[player3, thick, rounded corners=7pt] (-0.15, -2.25) rectangle (1.15, -1.75);
  \node[anchor=west, font=\footnotesize] at (0.04, -2) {$1-p$};
  \node[anchor=west, font=\footnotesize] at (1.04, -2) {$p$};
  \draw[->, thick] (0, -2.9) -- (0, -2.4); \node[font=\scriptsize] at (0, -3.05) {true state};
\end{tikzpicture}}

\begin{document}

\frame{\maketitle}

\begin{frame}{Where We Left Off}
Chapter 14: static games with incomplete information
\begin{itemize}
    \item Harsanyi transformation $\to$ a game with Nature
    \item Solution: BNE $=$ NE of that game
\end{itemize}
\bsk
Chapter 15: \textbf{dynamic} games with incomplete information
\begin{itemize}
    \item Same Harsanyi transformation
    \item But NE allows silly choices off the path (Chapter 4!)
    \item So use \textbf{WSE} (or SPNE) of the Harsanyi game
\end{itemize}
\bsk
Big applications: \textbf{reputation} and \textbf{signaling}
\vspace{-0.3in}
\end{frame}

% ==============================================================
% OPENING EXAMPLE (Bonanno Figs 15.1-15.5)
% ==============================================================

\begin{frame}{Example: Unsure of P2's Payoffs}
P1 knows her own payoffs. P1 thinks P2's payoffs are one of these:
\bsk
\begin{columns}[T]
\begin{column}{0.48\textwidth}
\centering
$\alpha$ (prob.\ $\tfrac13$)\\
\minigamealpha
\end{column}
\begin{column}{0.48\textwidth}
\centering
$\beta$ (prob.\ $\tfrac23$)\\
\minigamebeta
\end{column}
\end{columns}
\bsk
P2 knows which one. True state: $\alpha$
\bsk
Backward induction if P1 knew: at $\alpha$ play \underline{\hspace{0.3in}}. At $\beta$ play \underline{\hspace{0.3in}}
\vfill
{\footnotesize Bonanno Figures 15.1--15.3.}
\vspace{-0.2in}
\end{frame}

\begin{frame}{The Harsanyi Game}
\begin{columns}[c]
\begin{column}{0.5\textwidth}
\centering
\examplegame
\end{column}
\begin{column}{0.48\textwidth}
\centering
\begin{matrixgame}[top={DD, DA, AD, AA}, left={B, T}, player names={1, 2}, w=1.15]
  \payoffmatrix{1,0 & 1,0 & 1,0 & 1,0 \\ 0,\tfrac43 & \tfrac43,2 & \tfrac23,1 & 2,\tfrac53}
\end{matrixgame}\\[4pt]
{\footnotesize P2's strategy: (at $\alpha$, at $\beta$)}
\end{column}
\end{columns}
\bsk
\onslide<2->{BNE: \ \underline{\hspace{0.7in}} \ \underline{\hspace{0.7in}} \ \underline{\hspace{0.7in}}}
\bsk
\onslide<3->{Which ones have P2 doing something silly at a node? \underline{\hspace{1in}}}
\vspace{-0.3in}
\end{frame}

\begin{frame}{Use WSE Instead}
In dynamic games: BNE can rely on \textbf{incredible threats}
\begin{itemize}
    \item P2 is at a \textbf{singleton} node. P2 knows the state!
    \item $(B,DD)$: at $\beta$, P2 ``plans'' D: 1 instead of 2. Silly
    \item $(B,AD)$: silly at \textbf{both} nodes
\end{itemize}
\bsk
\textbf{From now on:} WSE (here it's the same as SPNE) \ $\Rightarrow$ \ \underline{\hspace{0.8in}}
\bsk
\onslide<2->{True state $\alpha$: what actually gets played? \underline{\hspace{0.8in}}\\
Backward induction in the true game says \underline{\hspace{0.3in}}}
\bsk
\onslide<3->{\textbf{Lesson:} P1's uncertainty can lead to play that looks ``wrong'' in the true game}
\vspace{-0.3in}
\end{frame}

% ==============================================================
% REPUTATION: THE CHAIN STORE
% ==============================================================

\begin{frame}[standout]
Reputation
\end{frame}

% ---- MobLab run: Centipede (replay) -----------------------------------
% Same settings as deck 05's run, deliberately replayed: students now
% know backward induction says Cash Out at once. Results come back in
% "Back to the Centipede", next to deck 05's numbers.
% ---------------------------------------------------------------------
\begin{frame}{MobLab: ``Centipede Game'', Again!}
\begin{itemize}
    \item MobLab time!
    \item Same game as Chapter 3: Invest or Cash Out
    \item You now know what backward induction says
    \item Play it the way you'd actually play it
\end{itemize}
\bsk
Write down: what did you expect your partner to do?
\vfill
\centering
Feel free to talk!
\vspace{-0.2in}
\end{frame}

\begin{frame}{Recall: The Chain Store (Chapter 3)}
IKEA at Polaris. Wayfair decides whether to open across I-71
\begin{itemize}
    \item Out: Wayfair 1, IKEA 5. \ Fight: (0, 0). \ Share: (2, 2)
\end{itemize}
\bsk
Then the same thing in Cincinnati, with another entrant
\bsk
IKEA's argument: ``Fight in Columbus. Cincinnati is watching!''
\begin{itemize}
    \item Fight, scare off Cincy: $0+5=5$ \quad Share, Cincy enters: $2+2=4$
\end{itemize}
\bsk
\onslide<2->{Backward induction: Cincy is the last city, so IKEA shares there no matter what. So fighting in Columbus buys nothing. \textbf{Both enter}}
\bsk
\onslide<3->{Deck 05: ``The reputation story needs \textbf{uncertainty}... Chapter 15''}
\vspace{-0.3in}
\end{frame}

\begin{frame}{Add a Little Doubt}
Two possible types of IKEA:
\begin{itemize}
    \item \textbf{Rational}: fight gives 0, share gives 2 (as before)
    \item \textbf{Hotheaded}: \textbf{likes} fighting. Fight gives 3, share gives 2
\end{itemize}
\bsk
\begin{columns}[c]
\begin{column}{0.58\textwidth}
\centering
\chainks
\end{column}
\begin{column}{0.4\textwidth}
\begin{itemize}
    \item IKEA really is rational
    \item Wayfair knows it
    \item Living Spaces (Cincy) isn't sure: hotheaded with prob.\ $p$
\end{itemize}
\end{column}
\end{columns}
\vfill
{\footnotesize Bonanno \S15.1 (Figures 15.6--15.11) uses slightly different numbers.}
\vspace{-0.2in}
\end{frame}

\begin{frame}{Solving It: Start in Cincinnati}
If Living Spaces enters Cincinnati:
\begin{itemize}
    \item Rational IKEA: \underline{\hspace{0.6in}} \qquad Hotheaded IKEA: \underline{\hspace{0.6in}}
\end{itemize}
\bsk
Living Spaces sees what happened in Columbus. Three possible histories:
\bsk
\centering
\renewcommand{\arraystretch}{1.6}
\begin{tabular}{|l|c|c|}
\hline
In Columbus... & Pr(hotheaded) & Living Spaces \\
\hline
Wayfair stayed Out & \hspace{0.8in} & \hspace{0.8in} \\
Wayfair In, IKEA Shared & & \\
Wayfair In, IKEA Fought & & \\
\hline
\end{tabular}
\bsk
\raggedright
Candidate: IKEA (both types) would fight in Columbus. Wayfair stays out
\vspace{-0.3in}
\end{frame}

\begin{frame}{Checking the WSE}
\textbf{Living Spaces after Out:} belief $=$ prior $p$ (Bayes)
$$\underbrace{p(0)+(1-p)(2)}_{\text{enter}}\ \le\ \underbrace{1\vphantom{()}}_{\text{stay out}} \iff p\ \ge\ \underline{\hspace{0.4in}}$$
\textbf{After Share / Fight:} off the path. WSE allows any beliefs. These are sensible
\bsk
\textbf{IKEA}, if Wayfair entered Columbus:
\begin{itemize}
    \item Rational: fight \underline{\hspace{0.6in}} vs.\ share \underline{\hspace{0.6in}}
    \item Hotheaded: fight \underline{\hspace{0.6in}} vs.\ share \underline{\hspace{0.6in}}
\end{itemize}
\textbf{Wayfair:} in \underline{\hspace{0.3in}} vs.\ out \underline{\hspace{0.3in}}
\vspace{-0.3in}
\end{frame}

\begin{frame}{Reputation Works!}
If Living Spaces thinks IKEA might be hotheaded ($p\ge\tfrac12$):
\begin{itemize}
    \item IKEA's threat to fight in Columbus is \textbf{credible}
    \item Even though Wayfair and IKEA both \textbf{know} IKEA is rational!
    \item IKEA fights to protect a reputation it doesn't deserve
\end{itemize}
\bsk
\onslide<2->{Compare Repeated Games: finite PD unravels with \textbf{complete} info}
\begin{itemize}
    \item<2-> A small chance of a ``crazy'' type can stop the unraveling
    \item<2-> Same idea explains passing in the centipede game
\end{itemize}
\bsk
\onslide<3->{It's also a \textbf{sequential equilibrium} (Bonanno Exercise 15.2)}
\vfill
{\footnotesize Not in the book: Kreps \& Wilson (1982); Milgrom \& Roberts (1982); Kreps, Milgrom, Roberts \& Wilson (1982).}
\vspace{-0.2in}
\end{frame}

\begin{frame}{Back to the Centipede}
\vspace{0.3in}
\centering
\renewcommand{\arraystretch}{2}
\begin{tabular}{l|c|c|}
\cline{2-3}
 & Chapter 3 (before BI) & Today (after BI) \\
\hline
\multicolumn{1}{|l|}{Avg.\ turns before Cash Out} & \hspace{1in} & \hspace{1in} \\
\hline
\multicolumn{1}{|l|}{\% Cash Out right away} & & \\
\hline
\end{tabular}
\bsk
\raggedright
If a few people \textbf{like} to invest, even selfish players should invest for a while
\vfill
{\footnotesize Not in the book. McKelvey \& Palfrey (1992, \textit{Econometrica}).}
\vspace{-0.2in}
\end{frame}

\begin{frame}{This Is Everywhere}
Building a reputation for toughness:
\begin{itemize}
    \item \textbf{Predatory pricing}: fight early entrants to scare later ones
    \item \textbf{Hostages}: ``we never negotiate with kidnappers''
    \item \textbf{Bosses}: fire the first slacker to keep the rest in line
    \item \textbf{Central banks}: fight inflation hard to be believed next time
    \item \textbf{Parents}: follow through on the threat... this time
\end{itemize}
\bsk
Real versions: many periods, many types, people learning slowly
\begin{itemize}
    \item Same core: act tough now so others' \textbf{beliefs} about you change
\end{itemize}
\vspace{-0.3in}
\end{frame}

% ==============================================================
% SIGNALING: STRIKES
% ==============================================================

\begin{frame}[standout]
Signaling
\end{frame}

\begin{frame}{Why Do Strikes Happen?}
A firm and a union negotiate wages. Then they strike. Then they agree
\begin{itemize}
    \item Strikes are costly for \textbf{both}: lost wages, lost production
    \item Why not just agree on day one?
\end{itemize}
\bsk
\onslide<2->{Bonanno's answer: \textbf{incomplete information}
\begin{itemize}
    \item The firm knows its profits. The union doesn't
    \item Enduring a strike is a \textbf{credible signal} that profits are low
\end{itemize}}
\vfill
{\footnotesize Bonanno \S15.1, Figures 15.12--15.15.}
\vspace{-0.2in}
\end{frame}

\begin{frame}{The Setup}
\begin{itemize}
    \item Union asks for a high wage $w_H=60$ or low wage $w_L=50$
    \item Firm says Yes (contract signed) or No
    \item No $\Rightarrow$ strike for one period. Then the union asks again
    \item No again $\Rightarrow$ no deal. Both get 0
    \item Waiting is costly: $\delta=0.6$
\end{itemize}
\bsk
Firm's profit (before wages): \textbf{High} $\pi_H=100$ (prob.\ $0.7$) or \textbf{Low} $\pi_L=55$
\begin{itemize}
    \item Firm knows. Union doesn't. Actually: Low
    \item Low firm can't afford the high wage: $55-60<0$
\end{itemize}
\bsk
\textbf{Last period is easy:} firm accepts any wage it can afford
\vspace{-0.3in}
\end{frame}

\begin{frame}{A Separating WSE}
\textbf{Candidate:} Union asks 60. High firm says Yes. Low firm says No (strike!)\\
Then union asks 50. Low firm says Yes
\bsk
\textbf{Union after a strike:} by Bayes, Pr(Low) $=$ \underline{\hspace{0.4in}} $\Rightarrow$ ask 50 \checkmark
\bsk
\textbf{High firm} won't fake being Low:
$$\underbrace{100-60}_{\text{pay 60 now}}\ \ \ge\ \ \underbrace{0.6\,(100-50)}_{\text{strike, then pay 50}} \qquad \underline{\hspace{0.4in}}\ \ge\ \underline{\hspace{0.4in}}\ \checkmark$$
\textbf{Low firm:} Yes to 60 gives \underline{\hspace{0.3in}}. \ No gives $0.6(55-50)=$ \underline{\hspace{0.3in}}
\bsk
\textbf{Union, first offer:} ask 50 $\to$ 50. \ Ask 60 $\to$ \ $0.7(60)+0.3(0.6)(50)=$ \underline{\hspace{0.4in}}
\vspace{-0.3in}
\end{frame}

\begin{frame}{What Happened?}
The true firm is Low $\Rightarrow$ \textbf{a strike actually happens}
\bsk
\begin{itemize}
    \item Inefficient: both lose from the delay
    \item But it's how the Low firm \textbf{proves} it's Low
    \item A High firm wouldn't sit through a strike just to save 10
\end{itemize}
\bsk
\textbf{Separating:} types choose differently, so actions reveal types
\bsk
\textbf{Signal:} works only if it's \textbf{too costly to fake} for the other type
\vspace{-0.3in}
\end{frame}

\begin{frame}{Signals Everywhere}
Costly actions that are cheaper for the ``good'' type:
\begin{itemize}
    \item \textbf{Education}: a degree is easier for productive workers (Spence)
    \item \textbf{Warranties}: cheap to offer if your product rarely breaks
    \item \textbf{Ads}: only worth paying for if customers come back
    \item \textbf{Engagement rings}: expensive on purpose
    \item \textbf{Peacocks}: a big tail is only survivable if you're strong
\end{itemize}
\bsk
Real versions: continuous signals, many types, noisy signals
\begin{itemize}
    \item Same core: the cost has to differ across types
\end{itemize}
\vfill
{\footnotesize Not in the book. Spence (1973, \textit{QJE}). 2001 Nobel Prize.}
\vspace{-0.2in}
\end{frame}

\begin{frame}{Truth in Advertising}
Gas station sign: ``Octane: at least 85.'' The truth is 89. True, but vague
\begin{itemize}
    \item Law: no false claims. But vague claims are legal
    \item Seller knows quality. Buyer doesn't
\end{itemize}
\bsk
\textbf{Skeptical buyer:} assume the \textbf{worst} quality consistent with the claim
\begin{itemize}
    \item ``At least 85'' $\Rightarrow$ ``it's 85''
    \item So a seller with 89 wants to say ``89''
    \item And then 88 does too, and 87...
\end{itemize}
\bsk
\textbf{WSE:} buyer ends up as if fully informed. No news is bad news!
\vfill
{\footnotesize Bonanno \S15.1 (Milgrom \& Roberts, 1986). See Exercise 15.4 for a naive buyer.}
\vspace{-0.2in}
\end{frame}

% ==============================================================
% WORKED EXAMPLE
% ==============================================================
% Answers are in the header comment.
% ---------------------------------------------------------------------

\begin{frame}[standout]
Let's Work One Out
\end{frame}

\begin{frame}{The Honors Section}
A senior is a \textbf{Star} (prob.\ $\tfrac13$) or \textbf{Average}. They know which
\begin{itemize}
    \item First: take a hard \textbf{Honors} capstone, or \textbf{Not}
    \item Honors costs a Star 1, an Average student 4 (more stress, worse grade)
\end{itemize}
\bsk
Nationwide sees the transcript (Honors or not), not the type
\begin{itemize}
    \item Hires as \textbf{Manager} or \textbf{Analyst}
    \item Student gets 5 as Manager, 2 as Analyst, minus Honors cost
    \item Nationwide: Manager pays off 4 with a Star, 0 with Average. Analyst: 2 either way
\end{itemize}
\vfill
{\footnotesize Honors doesn't make anyone better at the job. It's only a signal.}
\vspace{-0.2in}
\end{frame}

\begin{frame}{The Honors Section}
\begin{enumerate}
    \item Draw the Harsanyi tree. How many info sets does Nationwide have? \vspace{0.25in}
    \item For which beliefs does Nationwide hire a Manager? \vspace{0.25in}
    \item Find a \textbf{separating} WSE. Check every type and info set \vspace{0.35in}
    \item Find a \textbf{pooling} WSE where nobody takes Honors. What beliefs after Honors? \vspace{0.3in}
    \item Are those beliefs reasonable? Would an Average student ever take Honors? \vspace{0.2in}
    \item If Honors cost an Average student only 2, which WSE survive? \vspace{0.1in}
\end{enumerate}
\vspace{-0.3in}
\end{frame}

\begin{frame}{What Have We Learned, Charlie Brown?}
\begin{itemize}
    \item Dynamic + incomplete info: Harsanyi game, then \textbf{WSE}
    \begin{itemize}
        \item BNE alone allows incredible threats
    \end{itemize}
    \item \textbf{Reputation}: a little doubt about your type makes threats credible
    \begin{itemize}
        \item Fixes the chain-store paradox
    \end{itemize}
    \item \textbf{Signaling}: costly actions reveal types, if too costly to fake
    \begin{itemize}
        \item Strikes, education, warranties
    \end{itemize}
    \item Skeptical buyers + no lying $\Rightarrow$ full disclosure
\end{itemize}
\vfill
\vspace{-0.4in}
\end{frame}

\begin{frame}{Reading \& Practice}
\begin{itemize}
    \item Bonanno, Chapter 15
    \begin{itemize}
        \item \S15.1: read
        \item \S15.2 (bargaining, two-sided): skim
    \end{itemize}
    \item Exercises: \S15.3.1 (15.1, 15.3a, 15.5), \S15.3.2 (15.6)
    \item 15.2 and 15.3b need sequential equilibrium (Ch.~12). 15.7 is a challenge
    \item Solutions in \S15.4. Try them first
\end{itemize}
\vfill
\end{frame}

\end{document}
